% =============================================================================
%  PAES Studio: guía de síntesis (plantilla)
%  Lista para Overleaf: Nuevo proyecto → Subir proyecto → compilar con pdfLaTeX.
%
%  Formato: A4, dos columnas, cajas tcolorbox con la paleta Zinc de la app.
%  Entornos disponibles:
%    definicion, formula, ley, mecanismo, ideasclave, metodo, ejemplo,
%    errorfrecuente, dato, lineatiempo   (título opcional en [])
%  Comandos:
%    \alternativa{A}{texto}   alternativa con su tecla (igual que en la app)
%    \clave{C}        marca la alternativa correcta
%
%  El generador scripts/pipeline/generate_latex_summary.py reemplaza los bloques
%  delimitados por %%<NOMBRE>%% … %%</NOMBRE>%%. Fuera de ellos, edita libremente.
% =============================================================================
\documentclass[10pt,a4paper,twocolumn]{article}

\usepackage{lmodern}        % antes de fontenc/babel: no depende de las fuentes EC
\usepackage[T1]{fontenc}
\usepackage[utf8]{inputenc}
\usepackage[spanish,es-noshorthands,es-tabla]{babel}
\spanishplainpercent        % \% sin espaciado automático (evita choques en modo matemático)
\usepackage{textcomp}
\usepackage{amsmath,amssymb}
\usepackage[a4paper,margin=1.5cm,top=1.9cm,bottom=1.9cm,columnsep=0.75cm]{geometry}
\usepackage{xcolor}
\usepackage{enumitem}
\usepackage{booktabs}
\usepackage{array}
\usepackage{titlesec}
\usepackage{fancyhdr}
\usepackage[most]{tcolorbox}
\usepackage[hidelinks]{hyperref}

% ── Paleta Zinc (misma que la app) ───────────────────────────────────────────
\definecolor{zinc950}{HTML}{09090B}
\definecolor{zinc900}{HTML}{18181B}
\definecolor{zinc700}{HTML}{3F3F46}
\definecolor{zinc500}{HTML}{71717A}
\definecolor{zinc300}{HTML}{D4D4D8}
\definecolor{zinc200}{HTML}{E4E4E7}
\definecolor{zinc100}{HTML}{F4F4F5}
\definecolor{zinc50}{HTML}{FAFAFA}
\definecolor{esmeralda}{HTML}{047857}
\definecolor{rosa}{HTML}{BE123C}
\definecolor{ambar}{HTML}{B45309}

% ── Tipografía: texto con serifa, estructura en sans ─────────────────────────
\setlength{\parindent}{0pt}
\setlength{\parskip}{0.45em}
\linespread{1.08}
\color{zinc900}
\newcommand{\sen}{\operatorname{sen}}
\renewcommand{\arraystretch}{1.18}

\titleformat{\section}{\sffamily\bfseries\large\color{zinc950}}{\thesection}{0.6em}{}
  [{\color{zinc300}\titlerule[0.4pt]}]
\titleformat{\subsection}{\sffamily\bfseries\normalsize\color{zinc900}}{\thesubsection}{0.5em}{}
\titlespacing*{\section}{0pt}{1.1em}{0.5em}
\titlespacing*{\subsection}{0pt}{0.9em}{0.3em}

\setlist{leftmargin=1.2em,itemsep=0.15em,topsep=0.3em}
\setlist[itemize]{label={\color{zinc500}\textbullet}}

\pagestyle{fancy}
\fancyhf{}
\renewcommand{\headrulewidth}{0pt}
\fancyfoot[L]{\sffamily\scriptsize\color{zinc500}\guiaAsignatura}
\fancyfoot[R]{\sffamily\scriptsize\color{zinc500}\thepage}

% ── Cajas ────────────────────────────────────────────────────────────────────
\tcbset{
  paes/.style={
    enhanced, breakable, sharp corners=south, arc=1.5pt,
    boxrule=0.4pt, colback=white, colframe=zinc300,
    colbacktitle=white, coltitle=zinc900, titlerule=0pt,
    fonttitle=\sffamily\bfseries\footnotesize,
    toptitle=3pt, bottomtitle=0pt,
    left=7pt, right=7pt, top=2pt, bottom=6pt,
    before skip=7pt, after skip=7pt,
  },
  acento/.style={borderline west={2pt}{0pt}{#1}},
}
\newtcolorbox{definicion}[1][Definición]{paes, acento=zinc900, title={#1}}
\newtcolorbox{formula}[1][Fórmula clave]{paes, colback=zinc50, colbacktitle=zinc50, title={#1}}
\newtcolorbox{ejemplo}[1][Ejemplo resuelto]{paes, title={#1}}
\newtcolorbox{errorfrecuente}[1][Error frecuente]{paes, colframe=rosa!35, acento=rosa, title={#1}}
\newtcolorbox{dato}[1][Dato PAES]{paes, colframe=ambar!40, acento=ambar, title={#1}}
\newtcolorbox{mecanismo}[1][Mecanismo]{paes, colback=zinc50, colbacktitle=zinc50, acento=zinc700, title={#1}}
\newtcolorbox{ley}[1][Ley]{paes, colback=zinc50, colbacktitle=zinc50, acento=zinc900, title={#1}}
\newtcolorbox{ideasclave}[1][Ideas clave]{paes, title={#1}}
\newtcolorbox{metodo}[1][Cómo abordarlo en la PAES]{paes, acento=esmeralda, title={#1}}
% Línea de tiempo: filas «fecha & hecho \\».
\newenvironment{lineatiempo}[1][Línea de tiempo]
  {\begin{tcolorbox}[paes, acento=zinc700, title={#1}]\small\renewcommand{\arraystretch}{1.25}%
   \begin{tabular}{@{}>{\leavevmode\sffamily\bfseries\color{zinc700}}p{0.24\linewidth}p{\dimexpr0.76\linewidth-2\tabcolsep\relax}@{}}}
  {\end{tabular}\end{tcolorbox}}

% Alternativa con su tecla (eco de la anatomía A–E de la app).
\newtcbox{\tecla}{on line, boxrule=0.4pt, arc=1.5pt, colback=zinc50, colframe=zinc300,
  left=2.2pt, right=2.2pt, top=0.6pt, bottom=0.6pt, fontupper=\ttfamily\scriptsize\color{zinc700}}
\newtcbox{\teclaok}{on line, boxrule=0.4pt, arc=1.5pt, colback=esmeralda, colframe=esmeralda,
  left=2.2pt, right=2.2pt, top=0.6pt, bottom=0.6pt, fontupper=\ttfamily\scriptsize\bfseries\color{white}}
\newcommand{\alternativa}[2]{\par\noindent\hangindent=1.6em\hangafter=1\makebox[1.6em][l]{\tecla{#1}}#2\par}
\newcommand{\clave}[1]{\par\smallskip{\sffamily\footnotesize\color{esmeralda} Clave \teclaok{#1}}\par}

% Barra horizontal proporcional (para tablas de frecuencia).
\newcommand{\barra}[1]{{\color{zinc900}\rule[0.15ex]{#1\linewidth}{0.9ex}}}

%%<META>%%
\newcommand{\guiaTitulo}{Compendio de estudio}
\newcommand{\guiaAsignatura}{Competencia Matemática 1 (M1)}
\newcommand{\guiaDescripcion}{La prueba M1 evalúa la resolución de problemas, la modelación, la representación y la argumentación con los contenidos de 7.º básico a 2.º medio. Tiene 65 preguntas de 4 alternativas y dura 2 h 20 min.}
\newcommand{\guiaFuente}{4 ejes y 15 subtemas. Frecuencias calculadas sobre 7 ensayos (6 oficiales DEMRE), 390 preguntas. Generado el 1 de noviembre de 2026.}
%%</META>%%

\begin{document}

\twocolumn[{%
  \begin{tcolorbox}[enhanced, frame hidden, colback=white, left=0pt, right=0pt, top=0pt, bottom=6pt,
      borderline south={0.4pt}{0pt}{zinc300}]
    {\sffamily\footnotesize\color{zinc500} PAES Studio}\par\vspace{0.2em}
    {\sffamily\bfseries\LARGE\color{zinc950} \guiaTitulo: \guiaAsignatura}\par\vspace{0.35em}
    {\normalsize\color{zinc700} \guiaDescripcion}\par\vspace{0.25em}
    {\sffamily\scriptsize\color{zinc500} \guiaFuente}
  \end{tcolorbox}
  \vspace{0.6em}
}]

%%<CONTENIDO>%%
\section*{Cómo usar este compendio}

Cada subtema sigue el mismo orden: teoría, recuadro con las fórmulas, leyes o mecanismos centrales, ideas clave, método para responder en la PAES, trampas frecuentes y ejemplos resueltos paso a paso. Resuelve cada ejemplo antes de leer su desarrollo.

\section*{Qué se evalúa}

Distribución de las preguntas por eje en los ensayos disponibles en PAES Studio.

\begin{center}\small

\begin{tabular}{@{}p{0.42\linewidth}rrp{0.22\linewidth}@{}}

\toprule Eje & Preg. & \% & \\ \midrule

Números & 178 & 46 & \barra{0.100} \\

Probabilidad y estadística & 94 & 24 & \barra{0.053} \\

Geometría & 79 & 20 & \barra{0.045} \\

Álgebra y funciones & 39 & 10 & \barra{0.022} \\

\bottomrule

\end{tabular}

\end{center}

\begin{dato}Números y Probabilidad y estadística concentran el 70 \% de las preguntas: priorízalos si tienes poco tiempo.\end{dato}

\section{Números}

{\itshape\color{zinc700} Operatoria en los conjuntos numéricos, porcentajes, potencias, raíces y proporcionalidad.}

\subsection{Números enteros y racionales}

Los números enteros ($\mathbb{Z}$) incluyen los naturales, el cero y los negativos. Los racionales ($\mathbb{Q}$) son los que pueden escribirse como $\frac{a}{b}$ con $a$ y $b$ enteros y $b \neq 0$: todo decimal finito o periódico es racional.

Para operar se respeta la prioridad: primero paréntesis, luego potencias, después multiplicaciones y divisiones (de izquierda a derecha) y al final sumas y restas. Para sumar fracciones se iguala el denominador; para multiplicarlas se multiplican numeradores y denominadores; dividir es multiplicar por el recíproco.

En la recta numérica, un número es mayor mientras más a la derecha está. Sumar un negativo equivale a desplazarse a la izquierda; restar un negativo, a la derecha.

\begin{formula}[Operatoria con fracciones]
\begin{gather*}
\frac{a}{b}+\frac{c}{d}=\frac{ad+bc}{bd} \\
\frac{a}{b}\cdot\frac{c}{d}=\frac{ac}{bd} \\
\frac{a}{b}:\frac{c}{d}=\frac{a}{b}\cdot\frac{d}{c}
\end{gather*}

Decimal periódico: $0{,}\overline{3}=\dfrac{3}{9}=\dfrac{1}{3}$ y $0{,}1\overline{6}=\dfrac{16-1}{90}=\dfrac{1}{6}$.
\end{formula}

\begin{ideasclave}
\begin{itemize}
  \item Prioridad: paréntesis, potencias, $\times$ y $\div$, $+$ y $-$.
  \item Restar un negativo es sumar: $a-(-b)=a+b$.
  \item Todo decimal finito o periódico es racional.
  \item Para comparar fracciones, iguala denominadores o pásalas a decimal.
\end{itemize}
\end{ideasclave}

\begin{metodo}
\begin{enumerate}
  \item Identifica el orden de las operaciones antes de calcular.
  \item Simplifica las fracciones en cada paso para trabajar con números pequeños.
  \item Controla los signos: multiplica primero los signos y luego los valores absolutos.
  \item Estima el resultado (signo y tamaño) para descartar alternativas imposibles.
\end{enumerate}
\end{metodo}

\begin{errorfrecuente}[Trampas frecuentes del DEMRE]
\begin{itemize}
  \item Sumar numeradores y denominadores: $\frac{1}{2}+\frac{1}{3}\neq\frac{2}{5}$.
  \item Resolver de izquierda a derecha ignorando que la multiplicación va antes que la resta.
  \item Confundir $-3^2=-9$ con $(-3)^2=9$.
\end{itemize}
\end{errorfrecuente}

\begin{ejemplo}[Ejemplo resuelto]
{\sffamily\footnotesize\color{zinc500} Enunciado}\par
¿Cuál es el valor de $-2 + 3\cdot(-4) - (-5)$?
\alternativa{A}{$-19$}
\alternativa{B}{$-9$}
\alternativa{C}{$1$}
\alternativa{D}{$9$}
\clave{B}
{\sffamily\footnotesize\color{zinc500} Análisis}\par Primero la multiplicación; después sumas y restas, cuidando el doble signo.
{\sffamily\footnotesize\color{zinc500} Desarrollo}
\begin{enumerate}
  \item $3\cdot(-4)=-12$.
  \item $-(-5)=+5$.
  \item $-2-12+5=-9$.
\end{enumerate}
\end{ejemplo}

\begin{ejemplo}[Ejemplo tipo PAES 2: Operatoria con fracciones]
{\sffamily\footnotesize\color{zinc500} Enunciado}\par
¿Cuál es el valor de $\dfrac{2}{3} - \dfrac{1}{4}\cdot\dfrac{8}{9}$?
\alternativa{A}{$\dfrac{4}{9}$}
\alternativa{B}{$\dfrac{10}{27}$}
\alternativa{C}{$\dfrac{2}{9}$}
\alternativa{D}{$\dfrac{8}{9}$}
\clave{A}
{\sffamily\footnotesize\color{zinc500} Análisis}\par Respeta la prioridad de operaciones: primero la multiplicación y después la resta.
{\sffamily\footnotesize\color{zinc500} Desarrollo}
\begin{enumerate}
  \item Multiplica: $\dfrac{1}{4}\cdot\dfrac{8}{9} = \dfrac{8}{36} = \dfrac{2}{9}$.
  \item Escribe $\dfrac{2}{3}$ con denominador 9: $\dfrac{2}{3} = \dfrac{6}{9}$.
  \item Resta: $\dfrac{6}{9} - \dfrac{2}{9} = \dfrac{4}{9}$.
\end{enumerate}
\end{ejemplo}

\begin{ejemplo}[Ejemplo tipo PAES 3: Potencias]
{\sffamily\footnotesize\color{zinc500} Enunciado}\par
¿Cuál es el valor de $\dfrac{2^{5}\cdot 4^{3}}{8^{3}}$?
\alternativa{A}{$2$}
\alternativa{B}{$8$}
\alternativa{C}{$4$}
\alternativa{D}{$16$}
\clave{C}
{\sffamily\footnotesize\color{zinc500} Análisis}\par Escribe todas las potencias en la misma base (2) y opera con los exponentes.
{\sffamily\footnotesize\color{zinc500} Desarrollo}
\begin{enumerate}
  \item Como $4 = 2^2$ y $8 = 2^3$: $4^3 = (2^2)^3 = 2^6$ y $8^3 = (2^3)^3 = 2^9$.
  \item Multiplicación de igual base, se suman exponentes: $2^5\cdot 2^6 = 2^{11}$.
  \item División de igual base, se restan exponentes: $\dfrac{2^{11}}{2^{9}} = 2^{2} = 4$.
\end{enumerate}
\end{ejemplo}

\subsection{Porcentajes y problemas con dinero}

Un porcentaje es una razón con denominador 100: el $p\,\%$ de $x$ es $\frac{p}{100}\cdot x$. Aumentar un valor en un $p\,\%$ equivale a multiplicarlo por $\left(1+\frac{p}{100}\right)$, y disminuirlo, por $\left(1-\frac{p}{100}\right)$.

Los porcentajes sucesivos se multiplican, porque cada uno se aplica sobre el valor ya modificado. En problemas con dinero conviene ordenar los datos (precio unitario, cantidad, total) y revisar las unidades.

\begin{formula}[Factores de variación]
\begin{gather*}
\text{aumento de } p\,\%:\ x\left(1+\tfrac{p}{100}\right) \\
\text{descuento de } p\,\%:\ x\left(1-\tfrac{p}{100}\right)
\end{gather*}

\[\text{variación porcentual}=\frac{\text{valor final}-\text{valor inicial}}{\text{valor inicial}}\cdot 100\]
\end{formula}

\begin{ideasclave}
\begin{itemize}
  \item El 25 \% de 80 es $0{,}25\cdot 80=20$.
  \item Subir 10 \% y luego bajar 10 \% deja el valor un 1 \% más bajo: $1{,}1\cdot 0{,}9=0{,}99$.
  \item La variación porcentual siempre se calcula sobre el valor inicial.
  \item Precio con IVA (19 \%) $=1{,}19\cdot$ precio neto.
\end{itemize}
\end{ideasclave}

\begin{metodo}
\begin{enumerate}
  \item Identifica sobre qué cantidad (base) se aplica cada porcentaje.
  \item Convierte cada cambio en un factor multiplicativo.
  \item Multiplica los factores en orden y aplícalos a la base.
  \item Compara con las alternativas y verifica con un cálculo aproximado.
\end{enumerate}
\end{metodo}

\begin{errorfrecuente}[Trampas frecuentes del DEMRE]
\begin{itemize}
  \item Sumar porcentajes sucesivos: $+20\,\%$ y $-20\,\%$ no se anulan.
  \item Calcular la variación porcentual sobre el valor final.
  \item Confundir «el 30 \% de» con «un 30 \% más que».
\end{itemize}
\end{errorfrecuente}

\begin{ejemplo}[Ejemplo resuelto]
{\sffamily\footnotesize\color{zinc500} Enunciado}\par
Un producto cuesta \$12 000 sin IVA. Si el IVA es del 19 \%, ¿cuál es su precio final?
\alternativa{A}{\$2280}
\alternativa{B}{\$12 190}
\alternativa{C}{\$14 280}
\alternativa{D}{\$15 600}
\clave{C}
{\sffamily\footnotesize\color{zinc500} Análisis}\par El precio final es el neto más el 19 \% del neto, es decir, el neto multiplicado por 1,19.
{\sffamily\footnotesize\color{zinc500} Desarrollo}
\begin{enumerate}
  \item IVA: $0{,}19\cdot 12\,000=2280$.
  \item Precio final: $12\,000+2280=14\,280$, o directamente $12\,000\cdot 1{,}19$.
\end{enumerate}
\end{ejemplo}

\begin{ejemplo}[Ejemplo tipo PAES 2: Porcentajes sucesivos]
{\sffamily\footnotesize\color{zinc500} Enunciado}\par
Una zapatilla cuesta \$40 000. En marzo su precio sube un 20 \% y en abril el nuevo precio baja un 20 \%.

¿Cuál es el precio de la zapatilla después de ambos cambios?
\alternativa{A}{\$40 000}
\alternativa{B}{\$38 400}
\alternativa{C}{\$32 000}
\alternativa{D}{\$48 000}
\clave{B}
{\sffamily\footnotesize\color{zinc500} Análisis}\par Los porcentajes sucesivos se aplican sobre el precio vigente en cada momento, por eso se multiplican sus factores.
{\sffamily\footnotesize\color{zinc500} Desarrollo}
\begin{enumerate}
  \item Subir un 20 \% equivale a multiplicar por $1{,}2$: $40\,000\cdot 1{,}2 = 48\,000$.
  \item Bajar un 20 \% equivale a multiplicar por $0{,}8$: $48\,000\cdot 0{,}8 = 38\,400$.
  \item En una sola cuenta: $40\,000\cdot 1{,}2\cdot 0{,}8 = 40\,000\cdot 0{,}96 = 38\,400$. El precio final es un 4 \% menor que el inicial.
\end{enumerate}
\end{ejemplo}

\subsection{Potencias y raíces}

Una potencia $a^n$ abrevia la multiplicación de $a$ por sí mismo $n$ veces. Con exponentes enteros se cumple $a^0=1$ (si $a\neq 0$) y $a^{-n}=\frac{1}{a^n}$. La notación científica escribe un número como $k\cdot 10^n$ con $1\le k<10$.

La raíz $\sqrt[n]{a}$ es el número que elevado a $n$ da $a$; equivale a la potencia $a^{1/n}$. Solo se suman o restan raíces semejantes (mismo índice y misma cantidad subradical), por eso primero se simplifican extrayendo factores.

\begin{formula}[Propiedades]
\begin{gather*}
a^m\cdot a^n=a^{m+n} \\
\frac{a^m}{a^n}=a^{m-n} \\
(a^m)^n=a^{mn} \\
(ab)^n=a^nb^n
\end{gather*}

\begin{gather*}
\sqrt[n]{ab}=\sqrt[n]{a}\,\sqrt[n]{b} \\
\sqrt[n]{a^m}=a^{m/n} \\
\sqrt{72}=\sqrt{36\cdot 2}=6\sqrt{2}
\end{gather*}
\end{formula}

\begin{ideasclave}
\begin{itemize}
  \item Lleva todo a la misma base antes de operar: $8=2^3$, $27=3^3$.
  \item El exponente negativo invierte, no cambia el signo: $2^{-3}=\frac{1}{8}$.
  \item $\sqrt{a+b}\neq\sqrt{a}+\sqrt{b}$.
  \item Notación científica: $0{,}00034=3{,}4\cdot 10^{-4}$.
\end{itemize}
\end{ideasclave}

\begin{metodo}
\begin{enumerate}
  \item Descompón las bases en factores primos para igualarlas.
  \item Aplica una propiedad por paso y escríbelo.
  \item En raíces, extrae los factores cuadrados perfectos antes de sumar o restar.
  \item Comprueba el signo y el orden de magnitud del resultado.
\end{enumerate}
\end{metodo}

\begin{errorfrecuente}[Trampas frecuentes del DEMRE]
\begin{itemize}
  \item Multiplicar las bases: $2^3\cdot 2^4\neq 4^7$ (es $2^7$).
  \item Sumar exponentes al elevar una potencia: $(2^3)^2=2^6$, no $2^5$.
  \item Restar radicandos: $\sqrt{50}-\sqrt{18}\neq\sqrt{32}$.
\end{itemize}
\end{errorfrecuente}

\begin{ejemplo}[Ejemplo resuelto]
{\sffamily\footnotesize\color{zinc500} Enunciado}\par
¿Cuál es el valor de $\dfrac{9^{2}\cdot 3^{-1}}{27}$?
\alternativa{A}{$\dfrac{1}{3}$}
\alternativa{B}{$1$}
\alternativa{C}{$3$}
\alternativa{D}{$9$}
\clave{B}
{\sffamily\footnotesize\color{zinc500} Análisis}\par Todo se puede escribir en base 3.
{\sffamily\footnotesize\color{zinc500} Desarrollo}
\begin{enumerate}
  \item $9^2=(3^2)^2=3^4$ y $27=3^3$.
  \item Numerador: $3^4\cdot 3^{-1}=3^{3}$.
  \item $\dfrac{3^3}{3^3}=1$.
\end{enumerate}
\end{ejemplo}

\begin{ejemplo}[Ejemplo tipo PAES 2: Raíces]
{\sffamily\footnotesize\color{zinc500} Enunciado}\par
¿Cuál es el resultado de $\sqrt{50} - \sqrt{18}$?
\alternativa{A}{$\sqrt{32}$}
\alternativa{B}{$4$}
\alternativa{C}{$8\sqrt{2}$}
\alternativa{D}{$2\sqrt{2}$}
\clave{D}
{\sffamily\footnotesize\color{zinc500} Análisis}\par Solo se pueden restar raíces semejantes (mismo índice y misma cantidad subradical), así que primero se simplifican.
{\sffamily\footnotesize\color{zinc500} Desarrollo}
\begin{enumerate}
  \item Descompón: $\sqrt{50} = \sqrt{25\cdot 2} = 5\sqrt{2}$ y $\sqrt{18} = \sqrt{9\cdot 2} = 3\sqrt{2}$.
  \item Resta los coeficientes de las raíces semejantes: $5\sqrt{2} - 3\sqrt{2} = 2\sqrt{2}$.
\end{enumerate}
\end{ejemplo}

\subsection{Razones y proporcionalidad}

Dos magnitudes son directamente proporcionales si su cociente es constante ($\frac{y}{x}=k$): al duplicar una, la otra se duplica. Son inversamente proporcionales si su producto es constante ($xy=k$): al duplicar una, la otra se reduce a la mitad.

Una razón $a:b$ compara dos cantidades por cociente. Repartir una cantidad en la razón $a:b:c$ es dividirla en $a+b+c$ partes iguales y asignar $a$, $b$ y $c$ de esas partes.

\begin{formula}[Igualdades clave]
\begin{gather*}
\text{directa: } \frac{y_1}{x_1}=\frac{y_2}{x_2} \\
\text{inversa: } x_1\,y_1=x_2\,y_2
\end{gather*}
\end{formula}

{\sffamily\footnotesize\color{zinc500} Esquema comparativo}
\begin{center}\small
\begin{tabular}{@{}>{\raggedright\arraybackslash}p{\dimexpr(\linewidth-4\tabcolsep)/3\relax}>{\raggedright\arraybackslash}p{\dimexpr(\linewidth-4\tabcolsep)/3\relax}>{\raggedright\arraybackslash}p{\dimexpr(\linewidth-4\tabcolsep)/3\relax}@{}}
\toprule
 & \textbf{Directa} & \textbf{Inversa} \\
\midrule
Constante & $\frac{y}{x}=k$ & $x\cdot y=k$ \\
Gráfico & recta que pasa por el origen & hipérbola \\
Ejemplo & kilos comprados y precio & trabajadores y días de obra \\
\bottomrule
\end{tabular}
\end{center}

\begin{ideasclave}
\begin{itemize}
  \item Pregúntate: si una magnitud aumenta, ¿la otra aumenta o disminuye?
  \item Reparto proporcional: total $\div$ suma de las partes.
  \item Escala 1:50 000: 1 cm en el mapa son 50 000 cm = 500 m reales.
\end{itemize}
\end{ideasclave}

\begin{metodo}
\begin{enumerate}
  \item Decide con sentido común si la relación es directa o inversa.
  \item Plantea la igualdad de cocientes (directa) o de productos (inversa).
  \item Despeja la incógnita y revisa que el resultado tenga sentido.
\end{enumerate}
\end{metodo}

\begin{errorfrecuente}[Trampas frecuentes del DEMRE]
\begin{itemize}
  \item Tratar como directa una relación inversa (más trabajadores, menos días).
  \item En un reparto, dividir el total por el número de personas en vez de por la suma de las partes.
\end{itemize}
\end{errorfrecuente}

\begin{ejemplo}[Ejemplo resuelto]
{\sffamily\footnotesize\color{zinc500} Enunciado}\par
Se reparten \$60 000 entre Ana y Beto en la razón $2:3$. ¿Cuánto recibe Beto?
\alternativa{A}{\$24 000}
\alternativa{B}{\$30 000}
\alternativa{C}{\$36 000}
\alternativa{D}{\$40 000}
\clave{C}
{\sffamily\footnotesize\color{zinc500} Análisis}\par La razón $2:3$ divide el total en 5 partes iguales.
{\sffamily\footnotesize\color{zinc500} Desarrollo}
\begin{enumerate}
  \item Cada parte vale $\dfrac{60\,000}{2+3}=12\,000$.
  \item Beto recibe 3 partes: $3\cdot 12\,000=36\,000$.
\end{enumerate}
\end{ejemplo}

\begin{ejemplo}[Ejemplo tipo PAES 2: Proporcionalidad inversa]
{\sffamily\footnotesize\color{zinc500} Enunciado}\par
Si 6 personas pintan una casa en 10 días, ¿cuántos días tardarían 4 personas trabajando al mismo ritmo?
\alternativa{A}{$\dfrac{20}{3}$ días}
\alternativa{B}{12 días}
\alternativa{C}{15 días}
\alternativa{D}{20 días}
\clave{C}
{\sffamily\footnotesize\color{zinc500} Análisis}\par Menos personas tardan más días: es una proporcionalidad inversa, donde el producto se mantiene constante.
{\sffamily\footnotesize\color{zinc500} Desarrollo}
\begin{enumerate}
  \item El trabajo total es $6\cdot 10 = 60$ días-persona.
  \item Con 4 personas: $4\cdot d = 60 \Rightarrow d = 15$ días.
\end{enumerate}
\end{ejemplo}

\section{Álgebra y funciones}

{\itshape\color{zinc700} Expresiones algebraicas, ecuaciones, sistemas y funciones afín y cuadrática.}

\subsection{Expresiones algebraicas y factorización}

Una expresión algebraica combina números y letras con operaciones. Los términos semejantes (misma parte literal) se suman sumando sus coeficientes. Los productos notables permiten desarrollar rápido y, leídos al revés, factorizar.

Factorizar es escribir una expresión como producto. Primero se busca un factor común; después se reconoce un producto notable o un trinomio $x^2+bx+c$, que se factoriza buscando dos números cuyo producto sea $c$ y cuya suma sea $b$.

\begin{formula}[Productos notables]
\begin{gather*}
(a\pm b)^2=a^2\pm 2ab+b^2 \\
(a+b)(a-b)=a^2-b^2 \\
(x+p)(x+q)=x^2+(p+q)x+pq
\end{gather*}
\end{formula}

\begin{ideasclave}
\begin{itemize}
  \item Primero el factor común: $6x^2-9x=3x(2x-3)$.
  \item Diferencia de cuadrados: $x^2-49=(x+7)(x-7)$.
  \item $(a+b)^2\neq a^2+b^2$: falta el término $2ab$.
\end{itemize}
\end{ideasclave}

\begin{metodo}
\begin{enumerate}
  \item Busca un factor común numérico y literal.
  \item Cuenta los términos: dos sugieren diferencia de cuadrados; tres, cuadrado de binomio o trinomio $x^2+bx+c$.
  \item Comprueba multiplicando los factores obtenidos.
\end{enumerate}
\end{metodo}

\begin{errorfrecuente}[Trampas frecuentes del DEMRE]
\begin{itemize}
  \item Olvidar el doble producto en $(a+b)^2$.
  \item Equivocar los signos al sacar un factor negativo: $-x+3=-(x-3)$.
\end{itemize}
\end{errorfrecuente}

\begin{ejemplo}[Ejemplo resuelto]
{\sffamily\footnotesize\color{zinc500} Enunciado}\par
¿Cuál de las siguientes expresiones es equivalente a $4x^2-25$?
\alternativa{A}{$(2x-5)^2$}
\alternativa{B}{$(2x+5)(2x-5)$}
\alternativa{C}{$(4x-5)(x+5)$}
\alternativa{D}{$(2x-25)(2x+1)$}
\clave{B}
{\sffamily\footnotesize\color{zinc500} Análisis}\par Son dos cuadrados perfectos restados: $(2x)^2-5^2$.
{\sffamily\footnotesize\color{zinc500} Desarrollo}
\begin{enumerate}
  \item $4x^2=(2x)^2$ y $25=5^2$.
  \item Diferencia de cuadrados: $(2x)^2-5^2=(2x+5)(2x-5)$.
\end{enumerate}
\end{ejemplo}

\begin{ejemplo}[Ejemplo tipo PAES 2: Factorización]
{\sffamily\footnotesize\color{zinc500} Enunciado}\par
¿Cuál de las siguientes expresiones es equivalente a $x^2 - 9x + 20$?
\alternativa{A}{$(x-4)(x-5)$}
\alternativa{B}{$(x+4)(x+5)$}
\alternativa{C}{$(x-2)(x-10)$}
\alternativa{D}{$(x+2)(x-10)$}
\clave{A}
{\sffamily\footnotesize\color{zinc500} Análisis}\par Para factorizar $x^2 + bx + c$ se buscan dos números que multiplicados den $c$ y sumados den $b$.
{\sffamily\footnotesize\color{zinc500} Desarrollo}
\begin{enumerate}
  \item Se buscan dos números cuyo producto sea $20$ y cuya suma sea $-9$: son $-4$ y $-5$.
  \item Por lo tanto $x^2 - 9x + 20 = (x-4)(x-5)$.
  \item Comprobación: $(x-4)(x-5) = x^2 - 5x - 4x + 20 = x^2 - 9x + 20$.
\end{enumerate}
\end{ejemplo}

\subsection{Ecuaciones e inecuaciones lineales}

Una ecuación lineal se resuelve aplicando la misma operación a ambos lados hasta aislar la incógnita. En los problemas con enunciado lo central es traducir: «el doble de» es $2x$, «excede en 5» es $+5$, «la mitad de la suma» es $\frac{a+b}{2}$.

En una inecuación se procede igual, con una excepción: al multiplicar o dividir por un número negativo se invierte el sentido de la desigualdad. La solución es un intervalo.

\begin{formula}[Notación de intervalos (DEMRE)]
\begin{gather*}
[a,b]=\{x: a\le x\le b\} \\
]a,b[=\{x: a<x<b\} \\
[a,+\infty[=\{x: x\ge a\}
\end{gather*}
\end{formula}

\begin{ideasclave}
\begin{itemize}
  \item Si $-2x\le 6$, entonces $x\ge -3$ (se invierte la desigualdad).
  \item Para eliminar denominadores, multiplica \textbf{todos} los términos por el m.c.m.
  \item Revisa qué pide la pregunta: a veces no es $x$, sino otra cantidad.
\end{itemize}
\end{ideasclave}

\begin{metodo}
\begin{enumerate}
  \item Traduce el enunciado a una ecuación indicando qué representa la incógnita.
  \item Elimina paréntesis y denominadores.
  \item Agrupa términos y despeja; si divides por un negativo, invierte la desigualdad.
  \item Reemplaza la solución en el enunciado original para comprobar.
\end{enumerate}
\end{metodo}

\begin{errorfrecuente}[Trampas frecuentes del DEMRE]
\begin{itemize}
  \item No invertir la desigualdad al dividir por un número negativo.
  \item Multiplicar por el m.c.m. solo algunos términos.
  \item Responder el valor de $x$ cuando se pregunta por otra cantidad.
\end{itemize}
\end{errorfrecuente}

\begin{ejemplo}[Ejemplo resuelto]
{\sffamily\footnotesize\color{zinc500} Enunciado}\par
La suma de tres números enteros consecutivos es 72. ¿Cuál es el mayor de ellos?
\alternativa{A}{23}
\alternativa{B}{24}
\alternativa{C}{25}
\alternativa{D}{26}
\clave{C}
{\sffamily\footnotesize\color{zinc500} Análisis}\par Llama $x$ al menor; los otros son $x+1$ y $x+2$. Ojo: se pregunta por el mayor.
{\sffamily\footnotesize\color{zinc500} Desarrollo}
\begin{enumerate}
  \item $x+(x+1)+(x+2)=72 \Rightarrow 3x+3=72 \Rightarrow x=23$.
  \item Los números son 23, 24 y 25; el mayor es 25.
\end{enumerate}
\end{ejemplo}

\begin{ejemplo}[Ejemplo tipo PAES 2: Inecuaciones]
{\sffamily\footnotesize\color{zinc500} Enunciado}\par
¿Cuál es el conjunto solución de la inecuación $5 - 2x \le 11$?
\alternativa{A}{$x \le -3$}
\alternativa{B}{$x \ge -3$}
\alternativa{C}{$x \ge 3$}
\alternativa{D}{$x \le 3$}
\clave{B}
{\sffamily\footnotesize\color{zinc500} Análisis}\par Al multiplicar o dividir una desigualdad por un número negativo, el sentido de la desigualdad se invierte.
{\sffamily\footnotesize\color{zinc500} Desarrollo}
\begin{enumerate}
  \item Resta 5 a ambos lados: $-2x \le 6$.
  \item Divide por $-2$ e invierte el signo: $x \ge -3$.
  \item Comprobación con $x = 0$ (que cumple $x \ge -3$): $5 - 0 = 5 \le 11$. Verdadero.
\end{enumerate}
\end{ejemplo}

\subsection{Sistemas de ecuaciones lineales}

Un sistema de dos ecuaciones lineales con dos incógnitas representa dos rectas: tiene solución única si se cortan, ninguna si son paralelas e infinitas si son la misma recta.

Se resuelve por sustitución (despejar una incógnita y reemplazarla), igualación o reducción (amplificar y sumar o restar las ecuaciones para eliminar una incógnita).

\begin{formula}[Cantidad de soluciones de $ax+by=c,\textbackslash\{\} dx+ey=f$]
Única si $\dfrac{a}{d}\neq\dfrac{b}{e}$; ninguna si $\dfrac{a}{d}=\dfrac{b}{e}\neq\dfrac{c}{f}$; infinitas si $\dfrac{a}{d}=\dfrac{b}{e}=\dfrac{c}{f}$.
\end{formula}

\begin{ideasclave}
\begin{itemize}
  \item Reducción: elige la incógnita con coeficientes fáciles de igualar.
  \item Comprueba la solución en \textbf{ambas} ecuaciones.
  \item En problemas, define claramente qué representa cada variable.
\end{itemize}
\end{ideasclave}

\begin{metodo}
\begin{enumerate}
  \item Define las variables y plantea una ecuación por cada condición.
  \item Elimina una variable por reducción o sustitución.
  \item Calcula la otra variable y comprueba en las dos ecuaciones.
\end{enumerate}
\end{metodo}

\begin{errorfrecuente}[Trampas frecuentes del DEMRE]
\begin{itemize}
  \item Reemplazar el valor en la misma ecuación de la que se despejó (siempre se cumple).
  \item Responder la variable que no se pedía.
\end{itemize}
\end{errorfrecuente}

\begin{ejemplo}[Ejemplo resuelto]
{\sffamily\footnotesize\color{zinc500} Enunciado}\par
En una librería, 2 cuadernos y 3 lápices cuestan \$5400, y 1 cuaderno y 2 lápices cuestan \$3000. ¿Cuánto cuesta un cuaderno?
\alternativa{A}{\$600}
\alternativa{B}{\$1200}
\alternativa{C}{\$1800}
\alternativa{D}{\$2400}
\clave{C}
{\sffamily\footnotesize\color{zinc500} Análisis}\par Sean $c$ y $\ell$ los precios: $2c+3\ell=5400$ y $c+2\ell=3000$.
{\sffamily\footnotesize\color{zinc500} Desarrollo}
\begin{enumerate}
  \item De la segunda: $c=3000-2\ell$.
  \item Sustituye: $2(3000-2\ell)+3\ell=5400 \Rightarrow 6000-\ell=5400 \Rightarrow \ell=600$.
  \item $c=3000-1200=1800$. Comprobación: $2\cdot 1800+3\cdot 600=5400$.
\end{enumerate}
\end{ejemplo}

\begin{ejemplo}[Ejemplo tipo PAES 2: Ecuaciones lineales]
{\sffamily\footnotesize\color{zinc500} Enunciado}\par
¿Cuál es la solución de la ecuación $3(x-2) = 2x + 5$?
\alternativa{A}{$x = 3$}
\alternativa{B}{$x = 7$}
\alternativa{C}{$x = -1$}
\alternativa{D}{$x = 11$}
\clave{D}
{\sffamily\footnotesize\color{zinc500} Análisis}\par Distribuye el 3 en el paréntesis y luego agrupa los términos con $x$ en un lado.
{\sffamily\footnotesize\color{zinc500} Desarrollo}
\begin{enumerate}
  \item Distribuye: $3x - 6 = 2x + 5$.
  \item Resta $2x$ a ambos lados: $x - 6 = 5$.
  \item Suma 6: $x = 11$. Comprobación: $3(11-2) = 27$ y $2\cdot 11 + 5 = 27$.
\end{enumerate}
\end{ejemplo}

\begin{ejemplo}[Ejemplo tipo PAES 3: Sistemas de ecuaciones]
{\sffamily\footnotesize\color{zinc500} Enunciado}\par
En un taller de robótica hay 12 estudiantes. Si hay 4 mujeres más que hombres, ¿cuántas mujeres hay en el taller?
\alternativa{A}{4}
\alternativa{B}{6}
\alternativa{C}{8}
\alternativa{D}{10}
\clave{C}
{\sffamily\footnotesize\color{zinc500} Análisis}\par Dos condiciones, dos incógnitas: plantea un sistema y resuélvelo por reducción.
{\sffamily\footnotesize\color{zinc500} Desarrollo}
\begin{enumerate}
  \item Sean $m$ las mujeres y $h$ los hombres: $m + h = 12$ y $m - h = 4$.
  \item Suma ambas ecuaciones: $2m = 16 \Rightarrow m = 8$.
  \item Entonces $h = 12 - 8 = 4$; efectivamente hay 4 mujeres más que hombres.
\end{enumerate}
\end{ejemplo}

\subsection{Función lineal y afín}

Una función asigna a cada elemento del dominio un único valor. La función afín $f(x)=mx+n$ tiene como gráfico una recta: $m$ es la pendiente (cuánto cambia $f$ cuando $x$ aumenta en 1) y $n$ es el coeficiente de posición (corte con el eje $Y$). Si $n=0$ es lineal y modela una proporcionalidad directa.

En contexto, la pendiente es una tasa (pesos por kilómetro, litros por minuto) y $n$ es el valor inicial o cargo fijo.

\begin{formula}[Pendiente y ecuación]
\begin{gather*}
m=\frac{y_2-y_1}{x_2-x_1} \\
f(x)=m(x-x_1)+y_1
\end{gather*}
\end{formula}

\begin{ideasclave}
\begin{itemize}
  \item $m>0$: creciente; $m<0$: decreciente; $m=0$: constante.
  \item Evaluar $f(2)$ es reemplazar $x=2$; resolver $f(x)=2$ es buscar $x$.
  \item Cargo fijo + tarifa por unidad $\Rightarrow f(x)=\text{tarifa}\cdot x+\text{fijo}$.
\end{itemize}
\end{ideasclave}

\begin{metodo}
\begin{enumerate}
  \item Identifica el valor inicial (cuando la variable es 0) y la tasa de cambio.
  \item Calcula la pendiente con dos puntos, restando en el mismo orden.
  \item Escribe la función y verifícala con un punto conocido.
\end{enumerate}
\end{metodo}

\begin{errorfrecuente}[Trampas frecuentes del DEMRE]
\begin{itemize}
  \item Restar en distinto orden en el numerador y en el denominador de la pendiente.
  \item Confundir el cargo fijo con la tarifa por unidad.
\end{itemize}
\end{errorfrecuente}

\begin{ejemplo}[Ejemplo resuelto]
{\sffamily\footnotesize\color{zinc500} Enunciado}\par
Un taxi cobra \$400 de bajada de bandera y \$150 por cada 200 m recorridos. ¿Qué función entrega el cobro $C$, en pesos, de un viaje de $x$ metros?
\alternativa{A}{$C(x)=150x+400$}
\alternativa{B}{$C(x)=0{,}75x+400$}
\alternativa{C}{$C(x)=400x+150$}
\alternativa{D}{$C(x)=200x+150$}
\clave{B}
{\sffamily\footnotesize\color{zinc500} Análisis}\par La tarifa por metro es $\frac{150}{200}=0{,}75$ pesos; el cargo fijo es 400.
{\sffamily\footnotesize\color{zinc500} Desarrollo}
\begin{enumerate}
  \item Tarifa por metro: $150:200=0{,}75$.
  \item $C(x)=0{,}75x+400$. Control: 200 m cuestan $150+400=550$.
\end{enumerate}
\end{ejemplo}

\subsection{Función cuadrática}

La función $f(x)=ax^2+bx+c$ ($a\neq 0$) tiene como gráfico una parábola. Si $a>0$ abre hacia arriba y tiene un mínimo; si $a<0$ abre hacia abajo y tiene un máximo. Corta el eje $Y$ en $(0,c)$.

El vértice está en $x_v=-\frac{b}{2a}$ y la recta $x=x_v$ es el eje de simetría. Los ceros (cortes con el eje $X$) son las soluciones de $ax^2+bx+c=0$.

\begin{formula}[Vértice y ceros]
\begin{gather*}
x_v=-\frac{b}{2a},\quad y_v=f(x_v) \\
x=\frac{-b\pm\sqrt{b^2-4ac}}{2a}
\end{gather*}
\end{formula}

\begin{ideasclave}
\begin{itemize}
  \item El signo de $a$ decide si hay máximo o mínimo.
  \item El valor máximo o mínimo es la \textbf{ordenada} del vértice.
  \item Dos puntos simétricos tienen la misma altura: el vértice queda al medio.
\end{itemize}
\end{ideasclave}

\begin{metodo}
\begin{enumerate}
  \item Identifica $a$, $b$ y $c$ con su signo.
  \item Calcula $x_v=-\frac{b}{2a}$ y evalúa la función para obtener $y_v$.
  \item Lee qué pide el problema: ¿el instante (abscisa) o el valor máximo (ordenada)?
\end{enumerate}
\end{metodo}

\begin{errorfrecuente}[Trampas frecuentes del DEMRE]
\begin{itemize}
  \item Usar $\frac{b}{2a}$ sin el signo menos.
  \item Responder $x_v$ cuando se pide el valor máximo de la función.
\end{itemize}
\end{errorfrecuente}

\begin{ejemplo}[Ejemplo resuelto]
{\sffamily\footnotesize\color{zinc500} Enunciado}\par
La altura, en metros, de una pelota lanzada hacia arriba es $h(t)=-5t^2+20t$, con $t$ en segundos. ¿Cuál es la altura máxima?
\alternativa{A}{2 m}
\alternativa{B}{15 m}
\alternativa{C}{20 m}
\alternativa{D}{40 m}
\clave{C}
{\sffamily\footnotesize\color{zinc500} Análisis}\par $a=-5<0$: la parábola tiene un máximo en el vértice.
{\sffamily\footnotesize\color{zinc500} Desarrollo}
\begin{enumerate}
  \item $t_v=-\dfrac{20}{2\cdot(-5)}=2$ s.
  \item $h(2)=-5\cdot 4+20\cdot 2=20$ m.
\end{enumerate}
\end{ejemplo}

\begin{ejemplo}[Ejemplo tipo PAES 2: Función afín]
{\sffamily\footnotesize\color{zinc500} Enunciado}\par
El gráfico de la función $f(x) = mx + n$ pasa por los puntos $(1, 3)$ y $(3, 7)$. ¿Cuál es el valor de $f(5)$?
\alternativa{A}{$9$}
\alternativa{B}{$11$}
\alternativa{C}{$13$}
\alternativa{D}{$15$}
\clave{B}
{\sffamily\footnotesize\color{zinc500} Análisis}\par Con dos puntos se calcula la pendiente; con ella y un punto, el coeficiente de posición.
{\sffamily\footnotesize\color{zinc500} Desarrollo}
\begin{enumerate}
  \item Pendiente: $m = \dfrac{7-3}{3-1} = \dfrac{4}{2} = 2$.
  \item Usando $(1,3)$: $3 = 2\cdot 1 + n \Rightarrow n = 1$. Así, $f(x) = 2x + 1$.
  \item Evalúa: $f(5) = 2\cdot 5 + 1 = 11$.
\end{enumerate}
\end{ejemplo}

\begin{ejemplo}[Ejemplo tipo PAES 3: Función cuadrática]
{\sffamily\footnotesize\color{zinc500} Enunciado}\par
¿Cuáles son las coordenadas del vértice de la parábola $f(x) = x^2 - 4x + 3$?
\alternativa{A}{$(2, 3)$}
\alternativa{B}{$(4, 3)$}
\alternativa{C}{$(-2, 15)$}
\alternativa{D}{$(2, -1)$}
\clave{D}
{\sffamily\footnotesize\color{zinc500} Análisis}\par La abscisa del vértice de $ax^2 + bx + c$ es $x_v = -\dfrac{b}{2a}$; la ordenada se obtiene evaluando la función en ese valor.
{\sffamily\footnotesize\color{zinc500} Desarrollo}
\begin{enumerate}
  \item Aquí $a = 1$ y $b = -4$: $x_v = -\dfrac{-4}{2\cdot 1} = 2$.
  \item $f(2) = 2^2 - 4\cdot 2 + 3 = 4 - 8 + 3 = -1$.
  \item El vértice es $(2, -1)$. Como $a > 0$, es el punto mínimo de la parábola.
\end{enumerate}
\end{ejemplo}

\section{Geometría}

{\itshape\color{zinc700} Medición, teorema de Pitágoras, semejanza y transformaciones en el plano.}

\subsection{Perímetros, áreas y volúmenes}

El perímetro es la longitud del contorno; el área, la medida de la superficie; el volumen, el espacio que ocupa un cuerpo. Las figuras compuestas se descomponen en figuras conocidas, sumando o restando sus áreas.

Al multiplicar todas las longitudes por $k$, el perímetro se multiplica por $k$, el área por $k^2$ y el volumen por $k^3$.

\begin{formula}[Fórmulas]
Rectángulo: $A=b\cdot h$. Triángulo: $A=\dfrac{b\cdot h}{2}$. Círculo: $A=\pi r^2$, $P=2\pi r$.

Prisma y cilindro: $V=A_{\text{base}}\cdot h$ (cilindro: $V=\pi r^2h$). Cubo: $V=a^3$. Cono: $V=\dfrac{\pi r^2h}{3}$.
\end{formula}

\begin{ideasclave}
\begin{itemize}
  \item Descompón figuras compuestas antes de calcular.
  \item Revisa unidades: $1\ \text{m}^3=1000$ litros.
  \item Diámetro $=2\cdot$ radio.
\end{itemize}
\end{ideasclave}

\begin{metodo}
\begin{enumerate}
  \item Dibuja o marca las medidas conocidas.
  \item Elige la fórmula según la figura o cuerpo y verifica las unidades.
  \item Suma o resta las partes de la figura compuesta.
\end{enumerate}
\end{metodo}

\begin{errorfrecuente}[Trampas frecuentes del DEMRE]
\begin{itemize}
  \item Usar el diámetro como radio.
  \item Confundir área con perímetro.
  \item Olvidar dividir por 2 el área del triángulo.
\end{itemize}
\end{errorfrecuente}

\begin{ejemplo}[Ejemplo resuelto]
{\sffamily\footnotesize\color{zinc500} Enunciado}\par
Un estanque cilíndrico tiene 2 m de radio y 1,5 m de altura. Considerando $\pi\approx 3$, ¿cuál es su volumen?
\alternativa{A}{$9\ \text{m}^3$}
\alternativa{B}{$12\ \text{m}^3$}
\alternativa{C}{$18\ \text{m}^3$}
\alternativa{D}{$36\ \text{m}^3$}
\clave{C}
{\sffamily\footnotesize\color{zinc500} Análisis}\par Volumen del cilindro: área de la base por la altura.
{\sffamily\footnotesize\color{zinc500} Desarrollo}
\begin{enumerate}
  \item Base: $\pi r^2\approx 3\cdot 4=12\ \text{m}^2$.
  \item $V\approx 12\cdot 1{,}5=18\ \text{m}^3$.
\end{enumerate}
\end{ejemplo}

\begin{ejemplo}[Ejemplo tipo PAES 2: Áreas compuestas]
{\sffamily\footnotesize\color{zinc500} Enunciado}\par
A un rectángulo de 8 cm de largo y 5 cm de ancho se le recorta un semicírculo cuyo diámetro es uno de los lados de 8 cm. ¿Cuál es el área, en $\text{cm}^2$, de la figura que queda?
\alternativa{A}{$40 - 4\pi$}
\alternativa{B}{$40 - 16\pi$}
\alternativa{C}{$26 - 8\pi$}
\alternativa{D}{$40 - 8\pi$}
\clave{D}
{\sffamily\footnotesize\color{zinc500} Análisis}\par Área de la figura = área del rectángulo $-$ área del semicírculo.
{\sffamily\footnotesize\color{zinc500} Desarrollo}
\begin{enumerate}
  \item Rectángulo: $8\cdot 5 = 40\ \text{cm}^2$.
  \item El diámetro mide 8 cm, así que el radio es $r = 4$ cm (cabe dentro del ancho de 5 cm).
  \item Semicírculo: $\dfrac{\pi r^2}{2} = \dfrac{16\pi}{2} = 8\pi\ \text{cm}^2$.
  \item Área restante: $40 - 8\pi\ \text{cm}^2$.
\end{enumerate}
\end{ejemplo}

\subsection{Triángulos, Pitágoras y semejanza}

Los ángulos interiores de un triángulo suman $180^\circ$. En un triángulo rectángulo, el lado opuesto al ángulo recto es la hipotenusa y cumple el teorema de Pitágoras.

Dos triángulos son semejantes si tienen sus ángulos iguales; entonces sus lados correspondientes son proporcionales. El teorema de Thales permite calcular longitudes cuando hay rectas paralelas cortadas por secantes.

\begin{formula}[Relaciones clave]
\begin{gather*}
a^2+b^2=c^2 \\
\text{tríos: } (3,4,5),\ (5,12,13),\ (8,15,17)
\end{gather*}

Si la razón de semejanza es $k$, las áreas están en razón $k^2$.
\end{formula}

\begin{ideasclave}
\begin{itemize}
  \item La hipotenusa es el lado mayor.
  \item Los múltiplos de un trío pitagórico también lo son: (6, 8, 10).
  \item Sombras a la misma hora $\Rightarrow$ triángulos semejantes.
\end{itemize}
\end{ideasclave}

\begin{metodo}
\begin{enumerate}
  \item Identifica el ángulo recto y la hipotenusa.
  \item Plantea Pitágoras o la proporción entre lados correspondientes.
  \item Verifica que la hipotenusa resulte mayor que los catetos.
\end{enumerate}
\end{metodo}

\begin{errorfrecuente}[Trampas frecuentes del DEMRE]
\begin{itemize}
  \item Sumar los cuadrados cuando la incógnita es un cateto.
  \item Emparejar lados que no se corresponden en la semejanza.
\end{itemize}
\end{errorfrecuente}

\begin{ejemplo}[Ejemplo resuelto]
{\sffamily\footnotesize\color{zinc500} Enunciado}\par
Un rectángulo mide 9 cm de ancho y 12 cm de largo. ¿Cuánto mide su diagonal?
\alternativa{A}{13 cm}
\alternativa{B}{15 cm}
\alternativa{C}{21 cm}
\alternativa{D}{$\sqrt{63}$ cm}
\clave{B}
{\sffamily\footnotesize\color{zinc500} Análisis}\par La diagonal es la hipotenusa del triángulo rectángulo de catetos 9 y 12.
{\sffamily\footnotesize\color{zinc500} Desarrollo}
\begin{enumerate}
  \item $d^2=9^2+12^2=81+144=225$.
  \item $d=15$ cm (es el trío $3,4,5$ multiplicado por 3).
\end{enumerate}
\end{ejemplo}

\begin{ejemplo}[Ejemplo tipo PAES 2: Teorema de Pitágoras]
{\sffamily\footnotesize\color{zinc500} Enunciado}\par
Una escalera de 5 m de largo está apoyada en una pared vertical y su extremo superior alcanza una altura de 4 m. ¿A qué distancia de la pared está el pie de la escalera?
\alternativa{A}{$\sqrt{41}$ m}
\alternativa{B}{1 m}
\alternativa{C}{9 m}
\alternativa{D}{3 m}
\clave{D}
{\sffamily\footnotesize\color{zinc500} Análisis}\par Pared, suelo y escalera forman un triángulo rectángulo cuya hipotenusa es la escalera.
{\sffamily\footnotesize\color{zinc500} Desarrollo}
\begin{enumerate}
  \item Por Pitágoras: $d^2 + 4^2 = 5^2$.
  \item $d^2 = 25 - 16 = 9 \Rightarrow d = 3$ m.
\end{enumerate}
\end{ejemplo}

\begin{ejemplo}[Ejemplo tipo PAES 3: Semejanza]
{\sffamily\footnotesize\color{zinc500} Enunciado}\par
A cierta hora, una persona de 1,5 m de estatura proyecta una sombra de 2 m. A la misma hora, un árbol proyecta una sombra de 12 m. ¿Cuál es la altura del árbol?
\alternativa{A}{8 m}
\alternativa{B}{9 m}
\alternativa{C}{10,5 m}
\alternativa{D}{16 m}
\clave{B}
{\sffamily\footnotesize\color{zinc500} Análisis}\par Los rayos del sol llegan paralelos: persona y árbol forman triángulos semejantes, así que la razón altura/sombra es la misma.
{\sffamily\footnotesize\color{zinc500} Desarrollo}
\begin{enumerate}
  \item Plantea la proporción: $\dfrac{h}{12} = \dfrac{1{,}5}{2}$.
  \item Despeja: $h = 12\cdot\dfrac{1{,}5}{2} = 12\cdot 0{,}75 = 9$ m.
\end{enumerate}
\end{ejemplo}

\subsection{Plano cartesiano, vectores y transformaciones}

Las transformaciones isométricas (traslación, reflexión y rotación) cambian la posición de una figura sin alterar su forma ni su tamaño.

Trasladar un punto según un vector $\vec{v}=(a,b)$ es sumar las coordenadas. Las reflexiones y rotaciones con centro en el origen siguen reglas fijas sobre las coordenadas.

\begin{ley}[Reglas en coordenadas]
Traslación: $(x,y)\mapsto(x+a,\,y+b)$. Reflexión eje $X$: $(x,-y)$. Reflexión eje $Y$: $(-x,y)$.

Rotación en $90^\circ$ antihoraria: $(-y,x)$. Rotación en $180^\circ$: $(-x,-y)$.
\end{ley}

{\sffamily\footnotesize\color{zinc500} Esquema comparativo}
\begin{center}\small
\begin{tabular}{@{}>{\raggedright\arraybackslash}p{\dimexpr(\linewidth-4\tabcolsep)/3\relax}>{\raggedright\arraybackslash}p{\dimexpr(\linewidth-4\tabcolsep)/3\relax}>{\raggedright\arraybackslash}p{\dimexpr(\linewidth-4\tabcolsep)/3\relax}@{}}
\toprule
\textbf{Transformación} & \textbf{Regla} & \textbf{Conserva} \\
\midrule
Traslación por $(a,b)$ & $(x+a,\,y+b)$ & forma, tamaño y orientación \\
Reflexión eje $Y$ & $(-x,\,y)$ & forma y tamaño \\
Rotación $180^\circ$ & $(-x,\,-y)$ & forma y tamaño \\
\bottomrule
\end{tabular}
\end{center}

\begin{ideasclave}
\begin{itemize}
  \item Reflejar respecto del eje $Y$ cambia el signo de la abscisa.
  \item Una rotación de $180^\circ$ equivale a la simetría respecto del origen.
  \item Vector entre dos puntos: $\overrightarrow{AB}=B-A$.
\end{itemize}
\end{ideasclave}

\begin{metodo}
\begin{enumerate}
  \item Identifica el tipo de transformación y su elemento (vector, eje o centro).
  \item Aplica la regla de coordenadas.
  \item Verifica con un dibujo rápido en el plano.
\end{enumerate}
\end{metodo}

\begin{errorfrecuente}[Trampas frecuentes del DEMRE]
\begin{itemize}
  \item Confundir la reflexión en el eje $X$ con la del eje $Y$.
  \item Restar el vector de traslación en vez de sumarlo.
\end{itemize}
\end{errorfrecuente}

\begin{ejemplo}[Ejemplo resuelto]
{\sffamily\footnotesize\color{zinc500} Enunciado}\par
El punto $P(3,1)$ se rota en $90^\circ$ en sentido antihorario con centro en el origen. ¿Cuáles son las coordenadas de su imagen?
\alternativa{A}{$(1,-3)$}
\alternativa{B}{$(-1,3)$}
\alternativa{C}{$(-3,-1)$}
\alternativa{D}{$(3,-1)$}
\clave{B}
{\sffamily\footnotesize\color{zinc500} Análisis}\par Rotación de $90^\circ$ antihoraria: $(x,y)\mapsto(-y,x)$.
{\sffamily\footnotesize\color{zinc500} Desarrollo}
\begin{enumerate}
  \item $(3,1)\mapsto(-1,3)$.
  \item Comprobación: el punto pasa del primer al segundo cuadrante, como corresponde a un giro de un cuarto de vuelta antihorario.
\end{enumerate}
\end{ejemplo}

\begin{ejemplo}[Ejemplo tipo PAES 2: Transformaciones isométricas]
{\sffamily\footnotesize\color{zinc500} Enunciado}\par
¿Cuáles son las coordenadas de la imagen del punto $P(2, -3)$ al reflejarlo respecto del eje $Y$?
\alternativa{A}{$(-2, -3)$}
\alternativa{B}{$(2, 3)$}
\alternativa{C}{$(-2, 3)$}
\alternativa{D}{$(3, 2)$}
\clave{A}
{\sffamily\footnotesize\color{zinc500} Análisis}\par Una reflexión respecto del eje $Y$ cambia el signo de la abscisa y conserva la ordenada: $(x, y) \mapsto (-x, y)$.
{\sffamily\footnotesize\color{zinc500} Desarrollo}
\begin{enumerate}
  \item Aplica la regla a $P(2,-3)$: $(-2, -3)$.
  \item Interpretación: el punto queda a la misma distancia del eje $Y$, pero al otro lado.
\end{enumerate}
\end{ejemplo}

\begin{ejemplo}[Ejemplo tipo PAES 3: Vectores y traslaciones]
{\sffamily\footnotesize\color{zinc500} Enunciado}\par
El punto $A(1, 4)$ se traslada según el vector $\vec{v} = (-3, 2)$. ¿Cuáles son las coordenadas del punto trasladado?
\alternativa{A}{$(4, 2)$}
\alternativa{B}{$(-3, 8)$}
\alternativa{C}{$(-2, 6)$}
\alternativa{D}{$(2, -6)$}
\clave{C}
{\sffamily\footnotesize\color{zinc500} Análisis}\par Trasladar un punto según un vector es sumar sus coordenadas componente a componente.
{\sffamily\footnotesize\color{zinc500} Desarrollo}
\begin{enumerate}
  \item Suma las abscisas: $1 + (-3) = -2$ (el punto se mueve 3 unidades a la izquierda).
  \item Suma las ordenadas: $4 + 2 = 6$ (el punto sube 2 unidades).
  \item El punto trasladado es $A' = (-2, 6)$.
\end{enumerate}
\end{ejemplo}

\section{Probabilidad y estadística}

{\itshape\color{zinc700} Tablas y gráficos, medidas estadísticas y probabilidad.}

\subsection{Tablas, gráficos y medidas de tendencia central}

La media aritmética reparte el total en partes iguales; la mediana es el dato central de los datos ordenados; la moda, el dato más frecuente. La media es sensible a valores extremos; la mediana no.

En una tabla de frecuencias, la frecuencia absoluta cuenta cuántas veces aparece un dato; la relativa es esa cantidad dividida por el total, y la acumulada suma las frecuencias hasta ese dato.

\begin{formula}[Fórmulas]
\begin{gather*}
\bar{x}=\frac{\sum x_i\,f_i}{n} \\
\text{posición de la mediana}=\frac{n+1}{2} \\
f_r=\frac{f}{n}
\end{gather*}
\end{formula}

\begin{ideasclave}
\begin{itemize}
  \item Para la mediana, \textbf{ordena} primero los datos.
  \item Con $n$ par, la mediana es el promedio de los dos datos centrales.
  \item Un dato extremo mueve la media, pero casi no mueve la mediana.
\end{itemize}
\end{ideasclave}

\begin{metodo}
\begin{enumerate}
  \item Calcula el total de datos $n$ sumando las frecuencias.
  \item Usa la frecuencia acumulada para ubicar la posición de la mediana.
  \item Para la media, multiplica cada dato por su frecuencia antes de sumar.
\end{enumerate}
\end{metodo}

\begin{errorfrecuente}[Trampas frecuentes del DEMRE]
\begin{itemize}
  \item Calcular la mediana sin ordenar los datos.
  \item Promediar los valores de la tabla sin considerar sus frecuencias.
\end{itemize}
\end{errorfrecuente}

\begin{ejemplo}[Ejemplo resuelto]
{\sffamily\footnotesize\color{zinc500} Enunciado}\par
Las notas de un curso se resumen en la tabla:

\begin{center}\small
\begin{tabular*}{\linewidth}{@{\extracolsep{\fill}}lllll@{}}
\toprule
\textbf{Nota} & \textbf{4} & \textbf{5} & \textbf{6} & \textbf{7} \\
\midrule
Estudiantes & 3 & 5 & 8 & 4 \\
\bottomrule
\end{tabular*}
\end{center}

¿Cuál es la mediana de las notas?
\alternativa{A}{5}
\alternativa{B}{5,65}
\alternativa{C}{6}
\alternativa{D}{6,5}
\clave{C}
{\sffamily\footnotesize\color{zinc500} Análisis}\par Hay 20 datos: la mediana es el promedio de los datos en las posiciones 10 y 11.
{\sffamily\footnotesize\color{zinc500} Desarrollo}
\begin{enumerate}
  \item Frecuencias acumuladas: 3, 8, 16, 20.
  \item Las posiciones 10 y 11 caen en la nota 6 (posiciones 9 a 16).
  \item Mediana $=\dfrac{6+6}{2}=6$. (La media sería $\frac{113}{20}=5{,}65$.)
\end{enumerate}
\end{ejemplo}

\begin{ejemplo}[Ejemplo tipo PAES 2: Medidas de tendencia central]
{\sffamily\footnotesize\color{zinc500} Enunciado}\par
Los puntajes de un estudiante en cinco controles fueron:

\[4,\ 7,\ 7,\ 8,\ 9\]

Si agrega un sexto control con nota 13, ¿cuál es el nuevo promedio (media aritmética)?
\alternativa{A}{$7$}
\alternativa{B}{$7{,}5$}
\alternativa{C}{$8$}
\alternativa{D}{$8{,}5$}
\clave{C}
{\sffamily\footnotesize\color{zinc500} Análisis}\par La media es la suma de todos los datos dividida por la cantidad de datos.
{\sffamily\footnotesize\color{zinc500} Desarrollo}
\begin{enumerate}
  \item Suma original: $4 + 7 + 7 + 8 + 9 = 35$.
  \item Con el sexto dato: $35 + 13 = 48$.
  \item Nueva media: $\dfrac{48}{6} = 8$.
\end{enumerate}
\end{ejemplo}

\begin{ejemplo}[Ejemplo tipo PAES 3: Tablas de frecuencia]
{\sffamily\footnotesize\color{zinc500} Enunciado}\par
En un curso se preguntó a cada estudiante cuántos hermanos tiene. Los resultados se muestran en la tabla:

\begin{center}\small
\begin{tabular}{@{}ll@{}}
\toprule
\textbf{Número de hermanos} & \textbf{Frecuencia} \\
\midrule
0 & 5 \\
1 & 9 \\
2 & 4 \\
3 & 2 \\
\bottomrule
\end{tabular}
\end{center}

¿Qué porcentaje del curso tiene \textbf{al menos} 2 hermanos?
\alternativa{A}{20 \%}
\alternativa{B}{30 \%}
\alternativa{C}{45 \%}
\alternativa{D}{75 \%}
\clave{B}
{\sffamily\footnotesize\color{zinc500} Análisis}\par «Al menos 2» significa 2 o más. Suma esas frecuencias y compáralas con el total.
{\sffamily\footnotesize\color{zinc500} Desarrollo}
\begin{enumerate}
  \item Total de estudiantes: $5 + 9 + 4 + 2 = 20$.
  \item Con al menos 2 hermanos: $4 + 2 = 6$.
  \item Porcentaje: $\dfrac{6}{20}\cdot 100 = 30\,\%$.
\end{enumerate}
\end{ejemplo}

\subsection{Medidas de posición y dispersión}

Los cuartiles dividen los datos ordenados en cuatro partes con igual cantidad de datos: $Q_2$ es la mediana. El percentil $k$ es el valor bajo el cual queda aproximadamente el $k\,\%$ de los datos.

El rango (máximo menos mínimo) y el rango intercuartílico ($Q_3-Q_1$) miden la dispersión. El diagrama de caja muestra mínimo, $Q_1$, mediana, $Q_3$ y máximo.

\begin{definicion}[Definiciones]
Rango $=\text{máx}-\text{mín}$. Rango intercuartílico $=Q_3-Q_1$. Percentil 80: el 80 \% de los datos es menor o igual que él.
\end{definicion}

\begin{ideasclave}
\begin{itemize}
  \item Percentil alto no significa nota alta, sino posición relativa alta.
  \item En un diagrama de caja, la caja contiene el 50 \% central de los datos.
  \item Dos grupos pueden tener igual media y distinta dispersión.
\end{itemize}
\end{ideasclave}

\begin{metodo}
\begin{enumerate}
  \item Ordena los datos.
  \item Ubica la medida pedida por su posición (mediana, cuartil o percentil).
  \item Interpreta en el contexto: qué porcentaje de los datos queda por debajo.
\end{enumerate}
\end{metodo}

\begin{errorfrecuente}[Trampas frecuentes del DEMRE]
\begin{itemize}
  \item Interpretar el percentil 80 como «obtuvo el 80 \% del puntaje».
  \item Leer el largo de la caja como el rango total.
\end{itemize}
\end{errorfrecuente}

\begin{ejemplo}[Ejemplo resuelto]
{\sffamily\footnotesize\color{zinc500} Enunciado}\par
Ana obtuvo un puntaje ubicado en el percentil 80 de su curso. ¿Qué significa este resultado?
\alternativa{A}{Ana respondió correctamente el 80 \% de la prueba.}
\alternativa{B}{Ana obtuvo 80 puntos.}
\alternativa{C}{Aproximadamente el 80 \% del curso obtuvo un puntaje menor o igual que el de Ana.}
\alternativa{D}{Solo el 20 \% del curso aprobó la prueba.}
\clave{C}
{\sffamily\footnotesize\color{zinc500} Análisis}\par El percentil describe la posición de un dato respecto de los demás, no la cantidad de aciertos.
{\sffamily\footnotesize\color{zinc500} Desarrollo}
\begin{enumerate}
  \item Percentil 80: el 80 \% de los datos es menor o igual que ese valor.
  \item Por eso la alternativa C describe el resultado; A y B confunden posición con puntaje.
\end{enumerate}
\end{ejemplo}

\begin{ejemplo}[Ejemplo tipo PAES 2: Probabilidad clásica]
{\sffamily\footnotesize\color{zinc500} Enunciado}\par
Una bolsa contiene 3 bolitas rojas, 5 azules y 2 verdes, todas del mismo tamaño. Si se extrae una bolita al azar, ¿cuál es la probabilidad de que \textbf{no} sea azul?
\alternativa{A}{$\dfrac{1}{2}$}
\alternativa{B}{$\dfrac{3}{10}$}
\alternativa{C}{$\dfrac{1}{5}$}
\alternativa{D}{$\dfrac{2}{3}$}
\clave{A}
{\sffamily\footnotesize\color{zinc500} Análisis}\par Regla de Laplace: casos favorables dividido por casos posibles, cuando todos los resultados son equiprobables.
{\sffamily\footnotesize\color{zinc500} Desarrollo}
\begin{enumerate}
  \item Casos posibles: $3 + 5 + 2 = 10$ bolitas.
  \item Casos favorables (no azul): $3 + 2 = 5$.
  \item $P = \dfrac{5}{10} = \dfrac{1}{2}$. También: $P(\text{no azul}) = 1 - \dfrac{5}{10}$.
\end{enumerate}
\end{ejemplo}

\subsection{Probabilidad y conteo}

Si todos los resultados de un experimento son equiprobables, la probabilidad de un evento es la razón entre casos favorables y casos posibles (regla de Laplace). Toda probabilidad está entre 0 y 1.

Para contar casos se usa el principio multiplicativo o un diagrama de árbol. La probabilidad del complemento suele simplificar los problemas de «al menos uno».

\begin{formula}[Reglas]
\begin{gather*}
P(A)=\frac{\text{casos favorables}}{\text{casos posibles}} \\
P(A^c)=1-P(A)
\end{gather*}

\begin{gather*}
P(A\cup B)=P(A)+P(B)-P(A\cap B) \\
\text{independientes: } P(A\cap B)=P(A)\cdot P(B)
\end{gather*}
\end{formula}

\begin{ideasclave}
\begin{itemize}
  \item «Al menos uno» $=1-P(\text{ninguno})$.
  \item Con dos dados hay $6\cdot 6=36$ resultados ordenados.
  \item Eventos independientes: se multiplican sus probabilidades.
\end{itemize}
\end{ideasclave}

\begin{metodo}
\begin{enumerate}
  \item Describe el espacio muestral y verifica que sus resultados sean equiprobables.
  \item Cuenta los casos favorables (directamente o con el complemento).
  \item Simplifica la fracción y compárala con las alternativas.
\end{enumerate}
\end{metodo}

\begin{errorfrecuente}[Trampas frecuentes del DEMRE]
\begin{itemize}
  \item Contar pares no ordenados en experimentos con orden, como dos dados.
  \item Suponer equiprobables resultados que no lo son, como las sumas de dos dados.
\end{itemize}
\end{errorfrecuente}

\begin{ejemplo}[Ejemplo resuelto]
{\sffamily\footnotesize\color{zinc500} Enunciado}\par
Se lanza una moneda dos veces. ¿Cuál es la probabilidad de obtener al menos una cara?
\alternativa{A}{$\dfrac{1}{4}$}
\alternativa{B}{$\dfrac{1}{2}$}
\alternativa{C}{$\dfrac{2}{3}$}
\alternativa{D}{$\dfrac{3}{4}$}
\clave{D}
{\sffamily\footnotesize\color{zinc500} Análisis}\par Usa el complemento: «al menos una cara» es lo contrario de «ninguna cara».
{\sffamily\footnotesize\color{zinc500} Desarrollo}
\begin{enumerate}
  \item $P(\text{ninguna cara})=P(\text{SS})=\dfrac{1}{4}$.
  \item $P(\text{al menos una})=1-\dfrac{1}{4}=\dfrac{3}{4}$.
\end{enumerate}
\end{ejemplo}

\begin{ejemplo}[Ejemplo tipo PAES 2: Probabilidad con dos dados]
{\sffamily\footnotesize\color{zinc500} Enunciado}\par
Se lanzan dos dados comunes. ¿Cuál es la probabilidad de que la suma de los puntos obtenidos sea 7?
\alternativa{A}{$\dfrac{1}{6}$}
\alternativa{B}{$\dfrac{1}{12}$}
\alternativa{C}{$\dfrac{7}{36}$}
\alternativa{D}{$\dfrac{1}{11}$}
\clave{A}
{\sffamily\footnotesize\color{zinc500} Análisis}\par El espacio muestral tiene $6\cdot 6 = 36$ pares ordenados equiprobables.
{\sffamily\footnotesize\color{zinc500} Desarrollo}
\begin{enumerate}
  \item Pares con suma 7: $(1,6), (2,5), (3,4), (4,3), (5,2), (6,1)$, es decir, 6 casos.
  \item $P(\text{suma } 7) = \dfrac{6}{36} = \dfrac{1}{6}$.
\end{enumerate}
\end{ejemplo}
%%</CONTENIDO>%%

\end{document}
